\scanpage{37}\mspage{36}
aussitôt que $K$ est objet final, dans la catégorie de tous les carrés commutatifs
\diagram{\begin{tikzcd}[ampersand replacement=\&,column sep=large]
\unsure{{+}J'|J'}\arrow[r,"i'"]\arrow[d,"h_J"] \& I'\arrow[d,"h"]\\
J\arrow[r] \& I
\end{tikzcd}}
\deleted{qui en} \inserted{d'ens ordonnés (ou préord)} qui satisfont la condition précédente. Ça devrait impliquer $(b,c,d)\Longleftrightarrow a)$.

\uline{18 Juin}\quad Dans toutes ces questions, le point de vue des ultr. topologies \inserted{(i.e. top. \uuline{pré}ordonnées)} est peut-être le plus pratique, sans prendre la peine tjrs d'expliciter laborieusement en termes d'ensembles d'indices ordonnés. Une ``figure'' \inserted{de $L$} s'explicite comme un couple $(\deleted{A}S,T)$ d'une partie $S$ de $L$, et d'une ultr.topologie $T$ sur $S$ (topologie où \uline{toute} réunion de fermés est fermée, i.e \uline{toute} intersection d'ouverts est ouverte, i.e. l'intersection des voisinages \inserted{[\unsure{ouverts, si on veut|ouverts si on veut}]} de tout point est un voisinage, ou encore que tt pt a un plus petit voisinage (néc. ouvert). Si $(S,T),(S',T')$ sont comme ça, dire qu'ils sont compatibles signifie que $T,T'$ induisent la $=$ topologie sur $S\cap S'$, \inserted{qui est fermé dans $S'$ (pour $T$) et dans $S'$ (pour $T'$)} et alors y a une unique
